% Fichier généré automatiquement par tools/generate_corrections.py.
% Source : GEO/GEO-03/13_TransfoMouvement/13_TransfoMouvement.tex
% Statut : complet (5/5 question(s) renseignée(s), 0 reprise(s)).
\begin{corrigebox}[Corrigé extrait de la banque]
\CorrigeQuestion{1}
\begin{marginfigure}
\includegraphics[width=.4\linewidth]{13_01_c}
\end{marginfigure}

\CorrigeQuestion{2}
La fermeture géométrique s'écrit
\[
\vect{AB}+\vect{BC}+\vect{CA}=\vect{0},
\]
soit
\[
R\vect{i_1}-\lambda(t)\vect{i_2}+H\vect{j_0}=\vect{0}.
\]
La projection dans le repère $\rep{0}$ donne
\[
\left\{
\begin{array}{l}
R\cos\theta(t)-\lambda(t)\cos\varphi(t)=0,\\
R\sin\theta(t)-\lambda(t)\sin\varphi(t)+H=0.
\end{array}
\right.
\]
En élevant les deux relations au carré puis en les sommant,
\[
\lambda^2(t)=R^2+H^2+2HR\sin\theta(t).
\]
Comme $\lambda(t)$ représente une longueur,
\[
\boxed{\lambda(t)=\sqrt{R^2+H^2+2HR\sin\theta(t)}}.
\]

\CorrigeQuestion{3}
En dérivant,
\[
\boxed{
\dot{\lambda}(t)=
\frac{HR\dot{\theta}(t)\cos\theta(t)}
{\sqrt{R^2+H^2+2HR\sin\theta(t)}}
}.
\]

\CorrigeQuestion{4}
Pour un piston de section $S$,
\[
\boxed{q(t)=S\dot{\lambda}(t)}.
\]

\CorrigeQuestion{5}
Pour un piston de diamètre $D$,
\[
S=\frac{\pi D^2}{4}.
\]
À $60$ tours par minute,
\[
\omega=2\pi\ \mathrm{rad\,s^{-1}}
\quad\text{et}\quad
\theta(t)=\omega t.
\]
Le débit instantané est donc
\[
\boxed{
q(t)=
\frac{\pi D^2}{4}
\frac{HR\omega\cos(\omega t)}
{\sqrt{R^2+H^2+2HR\sin(\omega t)}}
}.
\]

\begin{minipage}[c]{.3\linewidth}
\includegraphics[width=\linewidth]{13_02_c}
\end{minipage}
\hfill
\begin{minipage}[c]{.3\linewidth}
\includegraphics[width=\linewidth]{13_03_c}
\end{minipage}
\hfill
\begin{minipage}[c]{.3\linewidth}
\includegraphics[width=\linewidth]{13_04_c}
\end{minipage}
\end{corrigebox}
